% v2.4.0 P1 figure batch 02: the final nine UIDs in the mandated order.
\documentclass[UTF8,a4paper,11pt,openany]{ctexbook}
\usepackage{../common/statlearnbook}
\SLSetBookMetadata{P1 绘图批次 02}{统计学习方法讲义 v2.4.0 绘图集成验证}{后九幅 P1 正式绘图源码局部编译}
\begin{document}
\pagestyle{empty}
\raggedbottom

\typeout{SLFIG-HARNESS-BEGIN FIG-P429-01}
% v2.3.1 figure UID: FIG-P429-01
% Proposed post-v2.3.1 polish for FIG-P429-01: protect key-label size and stroke hierarchy.
\tikzset{slfig-FIG-P429-01/.style={font=\fontsize{9.2pt}{11.0pt}\selectfont,every path/.append style={line cap=round,line join=round}}}
\begin{figure}[H]
\centering
\begin{tikzpicture}[slfig-FIG-P429-01,
  branch title/.style={font=\bfseries\fontsize{9.6pt}{11.4pt}\selectfont},
  branch note/.style={font=\fontsize{9pt}{10.6pt}\selectfont},
  source arrow/.style={-{Stealth[length=2.1mm,width=1.3mm]},line width=.92pt},
  local arrow/.style={-{Stealth[length=1.7mm,width=1mm]},line width=.68pt},
  evidence arrow/.style={-{Stealth[length=1.9mm,width=1.15mm]},line width=.75pt},
]
  \node[slfig node,minimum width=32mm,line width=.80pt,
        font=\fontsize{9.4pt}{11.2pt}\selectfont] (data) at (0,2.35) {观测数据 $X$};

  \begin{scope}[shift={(-4.25,0)}]
    \node[branch title,text=SLBlue] at (0,1.0) {聚类};
    \fill[SLBlue] (-.75,.20) circle (2pt);
    \fill[SLBlue] (-.50,-.12) circle (2pt);
    \fill[SLBlue] (-.92,-.30) circle (2pt);
    \node[draw=SLBlue,fill=white,circle,line width=.68pt,minimum size=4.2pt,inner sep=0pt] at (.55,.18) {};
    \node[draw=SLBlue,fill=white,circle,line width=.68pt,minimum size=4.2pt,inner sep=0pt] at (.82,-.17) {};
    \node[draw=SLBlue,fill=white,circle,line width=.68pt,minimum size=4.2pt,inner sep=0pt] at (.40,-.32) {};
    \draw[draw=SLBlue!72,line width=.62pt,densely dashed,rounded corners=2mm]
      (-1.15,.48) rectangle (-.25,-.55);
    \draw[draw=SLBlue!72,line width=.62pt,densely dashed,rounded corners=2mm]
      (.15,.48) rectangle (1.08,-.55);
    \node[branch note,text=SLBlue] at (0,-.88) {样本分组};
  \end{scope}

  \begin{scope}[shift={(0,0)}]
    \node[branch title,text=SLTeal] at (0,1.0) {降维};
    \foreach \x/\y in {-.85/.30,-.55/-.20,-.15/.38,.30/-.30,.78/.22}{\fill[SLTeal!72] (\x,\y) circle (1.7pt);}
    \draw[local arrow,draw=SLTeal]
      (-.85,.30)--(-.72,-.55);
    \draw[local arrow,draw=SLTeal]
      (-.15,.38)--(-.10,-.55);
    \draw[local arrow,draw=SLTeal]
      (.78,.22)--(.70,-.55);
    \draw[draw=SLTeal,line width=.96pt] (-1.1,-.55)--(1.1,-.55);
    \foreach \x in {-.72,-.10,.70}{\fill[SLTeal] (\x,-.55) circle (2pt);}
    \node[branch note,text=SLTeal] at (0,-.88) {高维点 $\to$ 低维坐标};
  \end{scope}

  \begin{scope}[shift={(4.25,0)}]
    \node[branch title,text=SLGold!88!black] at (0,1.0) {概率建模};
    \node[circle,draw=SLGold,fill=SLGoldBG,line width=.80pt,minimum size=7mm,inner sep=0pt] (z) at (0,.30) {$z$};
    \foreach \x/\name in {-0.80/o1,0/o2,0.80/o3}{
      \node[circle,draw=SLGold,fill=white,line width=.68pt,minimum size=5mm,inner sep=0pt] (\name) at (\x,-.45) {$x$};
      \draw[local arrow,draw=SLGold,line width=.82pt] (z)--(\name);
    }
    \node[branch note,text=SLGold!88!black] at (0,-.88) {潜变量生成观测};
  \end{scope}

  \draw[source arrow,draw=SLBlue] (data)--(-4.25,.88);
  \draw[source arrow,draw=SLTeal] (data)--(0,.88);
  \draw[source arrow,draw=SLGold] (data)--(4.25,.88);
  \node[slfig note,anchor=north,align=center,minimum width=46mm,line width=.72pt,
        font=\fontsize{9.2pt}{11pt}\selectfont] (evidence) at (0,-1.55)
    {可检验结构 / 评价证据};
  \draw[evidence arrow,draw=SLBlue] (-4.25,-.95)--(evidence.west);
  \draw[evidence arrow,draw=SLTeal] (0,-.95)--(evidence.north);
  \draw[evidence arrow,draw=SLGold] (4.25,-.95)--(evidence.east);
\end{tikzpicture}
\caption{聚类、降维和概率建模分别从样本分组、属性压缩和联合生成三个方向把观测映射为可检验的潜在结构。}\label{fig:V4-C01-three-structures}
\end{figure}

\clearpage
\typeout{SLFIG-HARNESS-END FIG-P429-01}

\typeout{SLFIG-HARNESS-BEGIN FIG-P578-01}
% v2.7.0 figure UID: FIG-P578-01
% One state machine; no whole-figure scaling; visible text has a 9.6pt source floor.
% Compact base-math runs use 10.7pt so their native 300 dpi ink clears 22px.
\pgfkeysifdefined{/tikz/alt}{}{\tikzset{alt/.code={}}}
\tikzset{slfig-FIG-P578-01/.style={font=\fontsize{9.6pt}{11.6pt}\selectfont,
  every path/.append style={line cap=round,line join=round}}}
\begin{figure}[H]
\newcommand{\slfigRejMath}[1]{{\fontsize{10.7pt}{12.4pt}\selectfont\ensuremath{#1}}}
\newcommand{\slfigRejLift}[1]{\phantom{#1}\llap{\smash{\raisebox{2.2pt}{#1}}}}
\centering
\begin{tikzpicture}[slfig-FIG-P578-01,>=Stealth,
  alt={对象—关系—结论：带预算的接受—拒绝状态机在任何随机调用前验证标量、密度、支持覆盖与全局包络证书，随后初始化空前缀和计数；每个成功候选使提议数加一，只有接受分支原子追加候选并增加接受数，普通拒绝保持前缀和接受数不变；每轮先检查接受数是否达到目标，再检查提议预算，因而同轮同时命中时返回 completed，随机源或数值失败保留最后合法前缀与失败位置，包络失败返回 invalid\_input 和 envelope\_condition\_failure 诊断。},
  every node/.style={font=\fontsize{9.6pt}{11.6pt}\selectfont,align=center},
  contract/.style={draw=SLRuleGray,rounded corners=2pt,fill=SLSoftGray,
    text width=75mm,minimum height=15mm,inner xsep=5pt,inner ysep=3pt,line width=.72pt},
  stepbox/.style={draw=SLRuleGray,rounded corners=2pt,fill=SLSoftGray,
    text width=45mm,minimum height=9mm,inner xsep=4pt,inner ysep=2pt,line width=.72pt},
  decision/.style={diamond,aspect=3.15,draw=SLBlue,fill=SLBlue!6,
    text width=35mm,minimum height=10mm,inner xsep=1.6pt,inner ysep=1.2pt,line width=.82pt},
  success/.style={draw=SLTeal,rounded corners=7pt,fill=SLTeal!9,
    text width=36mm,minimum height=11mm,inner xsep=3pt,inner ysep=2pt,line width=.88pt},
  budgetstop/.style={draw=SLGray,ellipse,densely dashed,fill=SLSoftGray,
    minimum width=37mm,minimum height=14mm,inner xsep=3pt,inner ysep=2pt,line width=.74pt},
  failure/.style={draw=SLRed,double,double distance=.9pt,fill=SLRedBG,
    text width=38mm,minimum height=11mm,inner xsep=3pt,inner ysep=3.3pt,line width=.72pt},
  ordinary/.style={draw=SLGray,densely dashed,rounded corners=2pt,fill=white,
    text width=34mm,minimum height=10mm,inner xsep=3pt,inner ysep=2pt,line width=.72pt},
  arr/.style={draw=SLBlue,line width=.84pt,-{Stealth[length=1.9mm]}},
  returnarr/.style={draw=SLGray,line width=.76pt,-{Stealth[length=1.8mm]}},
  edgeword/.style={font=\fontsize{9.6pt}{11.6pt}\selectfont,fill=white,inner sep=1pt},
  vedgeword/.style={edgeword,anchor=west,xshift=8mm,inner ysep=.3pt},
  vedgewordupper/.style={vedgeword,yshift=3.5pt},
]
  \node[contract] (precheck) at (0,9.50) {随机调用前完整验证：
    $N,B\in\mathbb N_0$，$1\le c<\infty$；$p,q$有限非负；
    $p>0\Rightarrow q>0$；$\rho=p/(cq)$在$q>0$处有限；\\
    $\operatorname{supp}(p)\subseteq\operatorname{supp}(q)$，且全局$p\le cq$};
  \node[decision] (valid) at (0,7.78) {全部输入与证书通过？};
  \node[failure] (badinput) at (6.15,7.23) {$\mathtt{invalid\_input}$\\
    异常出口；字段/位置\\\slfigRejMath{X_{1:a}=\varnothing,m=a=0}\\
    $\mathtt{envelope\_condition}$\\[-1pt]
    $\mathtt{\_failure}$（包络诊断）};

  \node[stepbox] (init) at (0,6.26) {初始化\slfigRejMath{m=a=0}，有\\
    \slfigRejLift{效样本前缀\slfigRejMath{X_{1:a}=\varnothing}}};
  \node[decision] (goal) at (0,4.85) {\slfigRejMath{a=N}？};
  \node[success] (completed) at (6.15,5.00) {$\mathtt{completed}$｜成功出口\\
    \slfigRejMath{X_{1:N},m,a=N}\\完成（含零目标）};
  \node[decision] (budgetcheck) at (0,3.15) {\slfigRejMath{m=B}？};
  \node[budgetstop] (budget) at (5.90,2.45) {$\mathtt{budget\_stop}$\\
    预算停止\\\slfigRejMath{X_{1:a},m=B,a<N}};

  \node[decision] (proposal) at (0,1.44) {调用\slfigRejMath{Y\sim q}成功？};
  \node[failure] (proposalfail) at (6.15,.09) {$\mathtt{random\_source\_failure}$\\
    \slfigRejMath{X_{1:a},m,a}\\位置\slfigRejMath{(m+1,\text{提议源})}};
  \node[stepbox] (countproposal) at (0,-.08) {候选生成成功：立即令\slfigRejMath{m\leftarrow m+1}};
  \node[decision] (uniform) at (0,-1.79) {均匀源成功并返回\slfigRejMath{U\sim\mathrm U(0,1)}？};
  \node[failure] (uniformfail) at (6.15,-2.59) {$\mathtt{random\_source\_failure}$\\
    \slfigRejMath{X_{1:a},m,a}\\位置\slfigRejMath{(m,\text{均匀源})}};

  \node[stepbox] (evaluate) at (0,-3.59) {求\slfigRejMath{p(Y),q(Y)}与\\
    \slfigRejLift{\slfigRejMath{\rho=p(Y)/(cq(Y))}}};
  \node[decision] (numericok) at (0,-5.29) {\slfigRejMath{p(Y),q(Y)}有限非负，\\\slfigRejMath{q(Y)>0}且\slfigRejMath{\rho}有限？};
  \node[failure] (numericfail) at (6.15,-5.39) {$\mathtt{numerical\_failure}$\\
    \slfigRejMath{X_{1:a},m,a}\\位置\slfigRejMath{(m,\text{密度/比率})}};
  \node[decision] (envelopeok) at (0,-7.29) {\slfigRejMath{0\le\rho\le1}？};
  \node[failure] (envelopefail) at (6.15,-8.19) {$\mathtt{invalid\_input}$\\
    \slfigRejMath{X_{1:a},m,a}；位置\slfigRejMath{m}\\
    $\mathtt{envelope\_condition}$\\[-1pt]
    $\mathtt{\_failure}$（包络诊断）};

  \node[decision] (accept) at (0,-9.02) {\slfigRejMath{U\le\rho}？};
  \node[stepbox,text width=34mm] (commit) at (-2.65,-10.45) {接受：原子追加\slfigRejMath{Y}，\\\slfigRejMath{a\leftarrow a+1}};
  \node[ordinary] (reject) at (2.65,-10.45) {普通拒绝：不追加，\\\slfigRejMath{a}与\slfigRejMath{X_{1:a}}不变};
  \coordinate (merge) at (0,-11.20);

  \draw[arr] (precheck)--(valid);
  \draw[arr] (valid.east)--node[edgeword,above]{否}(badinput.west);
  \draw[arr] (valid)--node[vedgewordupper]{是}(init);
  \draw[arr] (init)--(goal);
  \draw[arr] (goal.east)--node[edgeword,above]{是：优先}(completed.west);
  \draw[arr] (goal)--node[vedgeword]{否}(budgetcheck);
  \draw[arr] (budgetcheck.east)--node[edgeword,above]{是}(budget.west);
  \draw[arr] (budgetcheck)--node[vedgeword]{否}(proposal);
  \draw[arr] (proposal.east)--node[edgeword,above]{失败}(proposalfail.west);
  \draw[arr] (proposal)--node[vedgewordupper]{成功}(countproposal);
  \draw[arr] (countproposal)--(uniform);
  \draw[arr] (uniform.east)--node[edgeword,pos=.38,above]{失败}(uniformfail.west);
  \draw[arr] (uniform)--node[vedgewordupper]{成功}(evaluate);
  \draw[arr] (evaluate)--(numericok);
  \draw[arr] (numericok.east)--node[edgeword,above]{否}(numericfail.west);
  \draw[arr] (numericok)--node[vedgeword]{是}(envelopeok);
  \draw[arr] (envelopeok.east)--node[edgeword,above]{否}(envelopefail.west);
  \draw[arr] (envelopeok)--node[vedgeword]{是}(accept);
  \draw[arr] (accept.south west)--node[edgeword,left]{接受}(commit.north);
  \draw[arr] (accept.south east)--node[edgeword,right]{拒绝}(reject.north);
  \draw[returnarr] (commit.south)--(merge);
  \draw[returnarr] (reject.south)--(merge);
  \draw[returnarr] (merge)--(-4.75,-11.05)--
    node[edgeword,pos=.32,right,xshift=2pt]{先\slfigRejMath{a=N}\\再\slfigRejMath{m=B}}(-4.75,4.85)--(goal.west);
\end{tikzpicture}
\caption{带预算的接受--拒绝必须区分普通拒绝、收满目标、未收满而预算耗尽，以及包络证书失败。}\label{fig:V5-C02-rejection-flow}
\end{figure}

\clearpage
\typeout{SLFIG-HARNESS-END FIG-P578-01}

\typeout{SLFIG-HARNESS-BEGIN FIG-P605-01}
% v2.3.1 figure UID: FIG-P605-01
% v2.6.0 reverse redraw for FIG-P605-01: aligned semantic rows and collision-free spacing.
\tikzset{slfig-FIG-P605-01/.style={font=\fontsize{9.2pt}{11.0pt}\selectfont,every path/.append style={line cap=round,line join=round}}}
\begin{figure}[htbp]
\centering
\begin{tikzpicture}[slfig-FIG-P605-01,>=Stealth,every node/.style={font=\fontsize{9.2pt}{11.0pt}\selectfont,align=center},
  kernel/.style={draw=SLBlue,rounded corners=2pt,fill=white,minimum width=14mm,minimum height=9mm,line width=.72pt},
  choice/.style={diamond,aspect=2.1,draw=SLGold,fill=SLGold!8,text width=18mm,inner sep=1pt,line width=.78pt},
  flow/.style={-{Stealth[length=2.0mm]},draw=SLBlue,line width=.95pt},
  branch/.style={-{Stealth[length=1.8mm]},draw=SLGray,line width=.70pt}]
  \draw[draw=SLRuleGray,line width=.55pt,rounded corners=3pt,fill=SLSoftGray] (-7.4,-2.80) rectangle (-.55,2.55);
  \draw[draw=SLRuleGray,line width=.55pt,rounded corners=3pt,fill=SLSoftGray] (.55,-2.80) rectangle (7.4,2.55);
  \node[font=\fontsize{9.8pt}{11.8pt}\selectfont\bfseries,text=SLBlue] at (-3.98,2.02) {系统扫描：固定复合};
  \node[kernel] (k1) at (-6.20,.78) {$K_1$};
  \node[kernel] (k2) at (-4.40,.78) {$K_2$};
  \node (dots) at (-3.02,.78) {$\cdots$};
  \node[kernel] (kd) at (-1.55,.78) {$K_d$};
  \draw[flow] (k1)--(k2); \draw[flow] (k2)--(dots); \draw[flow] (dots)--(kd);
  \node at (-3.98,-.35) {$K_{\rm sys}=K_1K_2\cdots K_d$};
  \node[draw=SLGray,line width=.58pt,rounded corners=2pt,fill=white,minimum width=56mm,minimum height=7mm] at (-3.98,-2.05)
    {固定次序复合：通常不保证可逆};

  \node[font=\fontsize{9.8pt}{11.8pt}\selectfont\bfseries,text=SLBlue] at (3.98,2.02) {随机扫描：先抽坐标};
  \node[choice] (pick) at (3.98,1.05) {$J\sim\omega$};
  \node[kernel] (r1) at (1.55,-.30) {$K_1$};
  \node[kernel] (rj) at (3.98,-.30) {$K_j$};
  \node[kernel] (rd) at (6.35,-.30) {$K_d$};
  \draw[branch] (pick)--(r1); \draw[branch] (pick)--(rj); \draw[branch] (pick)--(rd);
  \node at (3.98,-1.20) {$K_{\rm rand}=\sum_{j=1}^{d}\omega_jK_j$};
  \node[draw=SLBlue,line width=.72pt,rounded corners=2pt,fill=white,
    text width=52mm,minimum height=10mm,inner xsep=3mm] at (3.98,-2.20)
    {若各 $K_j$ 关于 $\pi$ 可逆\\固定权重混合保持可逆};
\end{tikzpicture}
\caption{分量MH的系统扫描与随机扫描。固定顺序的核复合通常不保证可逆；可逆坐标核的固定权重随机混合仍保持可逆。}\label{fig:V5-C03-componentwise-sweep}
\end{figure}

\clearpage
\typeout{SLFIG-HARNESS-END FIG-P605-01}

\typeout{SLFIG-HARNESS-BEGIN FIG-P634-01}
% v2.7.0 figure UID: FIG-P634-01
% Structured glyph-safe redraw: preserve the one-arrow/eight-slot semantics,
% but express low-profile punctuation and equality through grouping and arrows.
\tikzset{
  slfig-FIG-P634-01/.style={font=\fontsize{9.6pt}{11.5pt}\selectfont,every path/.append style={line cap=round,line join=round}},
  sl634-old/.style={draw=SLRuleGray,densely dotted,fill=white,minimum width=13.2mm,minimum height=9.5mm,line width=.78pt,inner sep=2pt},
  sl634-done/.style={draw=SLBlue,fill=SLBlue!4,pattern=north east lines,pattern color=SLRuleGray,minimum width=13.2mm,minimum height=9.5mm,line width=.95pt,inner sep=2pt},
  sl634-halo/.style={draw=none,fill=white,inner xsep=1.5pt,inner ysep=1.25pt},
  sl634-now/.style={draw=SLGold,fill=SLGold!9,minimum width=13.2mm,minimum height=9.5mm,line width=1.05pt,inner sep=2pt},
  sl634-card/.style={draw=SLRuleGray,rounded corners=1.8pt,fill=white,line width=.55pt,inner xsep=5pt}
}
\pgfkeysifdefined{/tikz/alt}{}{\tikzset{alt/.code={}}}
\begin{figure}[htbp]
\centering
\begin{tikzpicture}[slfig-FIG-P634-01,>=Stealth,every node/.style={font=\fontsize{9.6pt}{11.5pt}\selectfont,align=center},
  alt={对象—关系—结论：固定顺序为 1,2,…,j−1,j,j+1,…,d；当前第 j 步的 x^[j] 前段是本轮新值且后段是前轮旧值；x^[d]=x^(t) 仅在末位更新结束后成立，并记录为轮末样本}]
  \node[font=\fontsize{10.6pt}{12.7pt}\selectfont\bfseries,text=SLBlue] at (0,2.16) {单轮系统扫描坐标带};
  \node at (-5.2,1.48) {$1$};
  \node at (-3.7,1.48) {$2$};
  \node at (-2.2,1.48) {省略};
  \node at (-.7,1.48) {前位};
  \node at (.8,1.48) {当前};
  \node at (2.3,1.48) {后位};
  \node at (3.8,1.48) {省略};
  \node at (5.3,1.48) {末位};
  \draw[-{Stealth[length=1.8mm,width=1.25mm]},draw=SLGray,line width=.60pt] (-5.52,1.10)--(5.62,1.10)
    node[right=4pt,font=\fontsize{9.6pt}{11.5pt}\selectfont] {更新顺序};
  \node[sl634-done] at (-5.2,0) {};
  \node[sl634-done] at (-3.7,0) {};
  \node[sl634-done] at (-2.2,0) {};
  \node[sl634-done] at (-.7,0) {};
  \node[sl634-halo] at (-5.2,0) {坐标\\$1$};
  \node[sl634-halo] at (-3.7,0) {坐标\\$2$};
  \node[sl634-halo] at (-2.2,0) {中间\\坐标};
  \node[sl634-halo] at (-.7,0) {前段\\末位};
  \node[sl634-now] at (.8,0) {当前\\坐标};
  \node[sl634-old] at (2.3,0) {后段\\首位};
  \node[sl634-old] at (3.8,0) {中间\\坐标};
  \node[sl634-old] at (5.3,0) {末位\\坐标};
  \node[font=\fontsize{9.6pt}{11.5pt}\selectfont,text=SLBlue] at (-2.9,-.76) {本轮新值};
  \node[font=\fontsize{9.6pt}{11.5pt}\selectfont,text=SLGold] at (.8,-.76) {当前新值};
  \node[font=\fontsize{9.6pt}{11.5pt}\selectfont,text=SLTextGray] at (3.8,-.76) {前轮旧值};

  \node[sl634-card,minimum width=118mm,minimum height=10.5mm] at (0,-1.52) {};
  \node[font=\fontsize{10.0pt}{12.0pt}\selectfont] at (0,-1.34) {$x^{[j]}$ 当前子步状态};
  \node[font=\fontsize{9.8pt}{11.7pt}\selectfont,text=SLBlue] at (-2.65,-1.81)
    {起始至当前坐标\quad 本轮新值};
  \node[font=\fontsize{9.8pt}{11.7pt}\selectfont,text=SLTextGray] at (2.65,-1.81)
    {后续至末位坐标\quad 前轮旧值};

  \node[sl634-card,minimum width=123mm,minimum height=11mm] at (0,-2.72) {};
  \node[inner sep=2.5pt,font=\fontsize{10.0pt}{12.0pt}\selectfont] (sweepend) at (-2.65,-2.87) {$x^{[d]}$};
  \node[inner sep=2.5pt,font=\fontsize{10.0pt}{12.0pt}\selectfont] (roundstate) at (.85,-2.87) {$x^{(t)}$};
  \node[inner sep=2.5pt,font=\fontsize{9.8pt}{11.7pt}\selectfont,text=SLTextGray] (sample) at (4.25,-2.87) {轮末样本};
  \draw[<->,draw=SLGray,line width=.60pt] (sweepend.east)--(roundstate.west)
    node[midway,above=4pt,font=\fontsize{9.6pt}{11.5pt}\selectfont,text=SLTextGray] {状态相同};
  \draw[-{Stealth[length=1.5mm,width=1.05mm]},draw=SLGray,line width=.60pt] (roundstate.east)--(sample.west)
    node[midway,above=4pt,font=\fontsize{9.6pt}{11.5pt}\selectfont,text=SLTextGray] {仅此记录};
\end{tikzpicture}
\captionsetup{width=.94\linewidth}
\caption[系统扫描逐坐标即时写回并区分子步状态与轮末样本]{系统扫描按固定次序即时写回；当前子步的前段使用本轮新值，后段沿用前轮旧值；末位更新结束后，末位状态与本轮样本状态相同并记录为轮末样本。}
\label{fig:V5-C04-coordinate-sweep}
\end{figure}

\clearpage
\typeout{SLFIG-HARNESS-END FIG-P634-01}

\typeout{SLFIG-HARNESS-BEGIN FIG-P657-01}
% v2.6.0 figure UID: FIG-P657-01
% Image-guided reverse redraw: preserve the six-node distribution grid, seven semantic relations and two legend samples.
\tikzset{slfig-FIG-P657-01/.style={font=\fontsize{9.2pt}{11.0pt}\selectfont,every path/.append style={line cap=round,line join=round}}}
\pgfkeysifdefined{/tikz/alt}{}{\tikzset{alt/.code={}}}
\begin{figure}[htbp]
\centering
\begin{tikzpicture}[slfig-FIG-P657-01,
  >=Stealth,
  every node/.style={font=\fontsize{9.4pt}{11.3pt}\selectfont,align=center},
  dist/.style={draw=SLTeal,fill=SLBlue!6,rounded corners=2pt,
    minimum width=30mm,minimum height=10.5mm,inner xsep=4pt,inner ysep=3pt,line width=.82pt},
  special/.style={draw=SLRuleGray,fill=SLSoftGray,rounded corners=2pt,
    minimum width=30mm,minimum height=10.5mm,inner xsep=4pt,inner ysep=3pt,line width=.78pt},
  contain/.style={-{Stealth[open,length=1.9mm,width=1.4mm]},line width=.58pt,draw=SLTextGray},
  conj/.style={-{Stealth[length=2.15mm,width=1.55mm]},line width=1pt,draw=SLDeepBlue},
  edgelabel/.style={font=\fontsize{8.8pt}{10.5pt}\selectfont,text=SLTextGray,inner sep=0pt},
  alt={对象—关系—结论：六个常用分布之间同时存在特殊情形关系和共轭关系：Beta是二维Dirichlet，二项是二维多项，类别分布是单次多项，Bernoulli同时是单次二项；粗箭头表示共轭先验而不是集合包含}]
\node[font=\fontsize{9.5pt}{11.4pt}\selectfont\bfseries,text=SLBlue] at (-5.30,2.15) {先验族};
\node[dist] (dir) at (-2.35,2.15) {Dirichlet分布};
\node[dist] (beta) at (2.80,2.15) {Beta分布\\$K=2$};
\node[font=\fontsize{9.5pt}{11.4pt}\selectfont\bfseries,text=SLBlue] at (-5.30,0) {似然族};
\node[dist] (multi) at (-2.35,0) {多项分布};
\node[dist] (binom) at (2.80,0) {二项分布\\$K=2$};
\node[font=\fontsize{9.5pt}{11.4pt}\selectfont\bfseries,text=SLBlue] at (-5.30,-2.15) {单次试验};
\node[special] (cat) at (-2.35,-2.15) {类别分布\\$N=1$};
\node[special] (bern) at (2.80,-2.15) {Bernoulli分布\\$K=2,N=1$};

\draw[conj] (dir)--(multi);
\draw[conj] (beta)--(binom);
\draw[contain] (dir)--node[edgelabel,above=5.4pt]{特殊情形}(beta);
\draw[contain] (multi)--node[edgelabel,above=5.4pt]{特殊情形}(binom);
\draw[contain] (multi)--node[edgelabel,pos=.52,right=6pt]{$N=1$}(cat);
\draw[contain] (binom)--node[edgelabel,pos=.52,right=6pt]{$N=1$}(bern);
\draw[contain] (cat)--node[edgelabel,above=5.4pt]{$K=2$}(bern);
\draw[conj] (4.90,1.15)--(6.15,1.15)
  node[edgelabel,right=5.5pt]{共轭};
\draw[contain] (4.90,.30)--(6.15,.30)
  node[edgelabel,right=5.5pt]{特殊情形};
\end{tikzpicture}
\captionsetup{width=.94\linewidth}
\caption{六个常用分布之间同时存在特殊情形关系和共轭关系：Beta是二维Dirichlet，二项是二维多项，类别分布是单次多项，Bernoulli同时是单次二项；粗箭头表示共轭先验而不是集合包含}
\label{fig:V5-C05-distribution-relations}
\end{figure}

\clearpage
\typeout{SLFIG-HARNESS-END FIG-P657-01}

\typeout{SLFIG-HARNESS-BEGIN FIG-P683-01}
% v2.6.0 figure UID: FIG-P683-01
% Semantic lock: 18 node / 5 draw, N_m nested in M, K separate; no resizebox.
\tikzset{slfig-FIG-P683-01/.style={font=\fontsize{9.4pt}{11.3pt}\selectfont,every path/.append style={line cap=round,line join=round}}}
\pgfkeysifdefined{/tikz/alt}{}{\tikzset{alt/.code={}}}
\begin{figure}[htbp]
\centering
\begin{tikzpicture}[slfig-FIG-P683-01,slfig high,>=Stealth,
  every node/.style={font=\fontsize{9.4pt}{11.3pt}\selectfont},
  latent/.style={circle,draw=SLTeal,fill=SLTealBG,line width=.80pt,minimum size=8.2mm,inner sep=0pt,font=\fontsize{9.6pt}{11.5pt}\selectfont},
  hyper/.style={draw=none,fill=none,inner sep=1pt,text=SLGray,font=\fontsize{9.6pt}{11.5pt}\selectfont},
  obs/.style={circle,draw=SLBlue,fill=SLBlue,line width=.85pt,text=white,minimum size=8.2mm,inner sep=0pt,font=\bfseries\fontsize{9.6pt}{11.5pt}\selectfont},
  edge/.style={-{Stealth[length=2.1mm,width=1.5mm]},draw=SLBlue,line width=.95pt},
  plate/.style={draw=SLGray!62,line width=.66pt,rounded corners=2.2pt,dash pattern=on 2.6pt off 1.8pt},
  plabel/.style={text=SLGray,font=\fontsize{9.2pt}{11.0pt}\selectfont,inner sep=0pt},
  legendtext/.style={anchor=west,font=\fontsize{9.2pt}{11.0pt}\selectfont,inner sep=0pt},
  alt={对象—关系—结论：完整Bayes LDA盘式图把超参数、潜变量和观测变量分开：每篇文档共享一个主题比例，每个词位拥有一个主题指派，所有文档共享带Dirichlet先验的主题词分布；盘框标明重复次数，箭头只表示条件依赖方向}]
\node[hyper] (alpha) at (-5.95,1.25) {$\alpha$};
\node[latent] (theta) at (-3.05,1.25) {$\theta_m$};
\node[latent] (z) at (-.65,1.25) {$z_{mn}$};
\node[obs] (w) at (1.65,1.25) {$w_{mn}$};
\node[hyper] (beta) at (-5.95,-2.10) {$\beta$};
\node[latent] (phi) at (-.65,-2.10) {$\varphi_k$};

\node[plate,fit=(z)(w),inner xsep=4.6mm,inner ysep=4.4mm] (nplate) {};
\node[plate,fit=(theta)(nplate),inner xsep=5.2mm,inner ysep=6.8mm] (mplate) {};
\node[plate,fit=(phi),inner xsep=4.5mm,inner ysep=3.5mm] (kplate) {};

\draw[edge] (alpha) -- (theta);
\draw[edge] (theta) -- (z);
\draw[edge] (z) -- (w);
\draw[edge] (beta) -- (phi);
\draw[edge] (phi) -- (w);

\node[plabel,anchor=south east] at ($(nplate.north east)+(-1.5mm,1.55mm)$) {$N_m$ 个词位};
\node[plabel,anchor=north east] at ($(mplate.south east)+(0,-1.9mm)$) {$M$ 篇文档};
\node[plabel,anchor=north east] at ($(kplate.south east)+(0,-1.9mm)$) {$K$ 个主题};
\node[obs,minimum size=6.4mm] at (5.45,.95) {};
\node[legendtext] at (6.25,.95) {观测变量};
\node[latent,minimum size=6.4mm] at (5.45,0) {};
\node[legendtext] at (6.25,0) {潜变量};
\node[hyper] at (5.45,-.95) {$\alpha,\beta$};
\node[legendtext] at (6.25,-.95) {超参数（plate 外）};
\end{tikzpicture}
\captionsetup{width=.94\linewidth}
\caption{完整Bayes LDA盘式图把超参数、潜变量和观测变量分开：每篇文档共享一个主题比例，每个词位拥有一个主题指派，所有文档共享带Dirichlet先验的主题词分布；盘框标明重复次数，箭头只表示条件依赖方向}
\label{fig:V5-C06-plate-graph}
\end{figure}

\clearpage
\typeout{SLFIG-HARNESS-END FIG-P683-01}

\typeout{SLFIG-HARNESS-BEGIN FIG-P694-01}
% v2.7.0 figure UID: FIG-P694-01
% R2.4 compact matrix/positioning layout; no overall scaling.
\tikzset{slfig-FIG-P694-01/.style={
  font=\fontsize{9.6pt}{10.3pt}\selectfont,
  every path/.append style={line cap=round,line join=round}}}
\pgfkeysifdefined{/tikz/alt}{}{\tikzset{alt/.code={}}}
\begin{figure}[htbp]
\centering
\makebox[\linewidth][c]{%
\begin{tikzpicture}[slfig-FIG-P694-01,
  every node/.style={
    font=\fontsize{9.6pt}{10.3pt}\selectfont,
    align=center},
  box/.style={
    draw=SLDeepBlue,fill=white,rounded corners=2pt,
    line width=.78pt,inner xsep=1.6pt,inner ysep=1.6pt},
  localbox/.style={
    box,draw=SLMidBlue,fill=SLSoftTeal},
  gate/.style={
    box,double,double distance=.55pt,line width=.76pt},
  localgate/.style={
    gate,draw=SLMidBlue,fill=SLSoftTeal},
  cert/.style={
    gate,fill=SLSoftBlue},
  localcert/.style={
    localgate},
  ok/.style={
    draw=SLSelfCheckTeal,fill=SLSoftTeal,rounded corners=2pt,
    line width=.82pt,inner xsep=1.6pt,inner ysep=1.6pt},
  budget/.style={
    draw=SLAccentOrange,fill=SLWarmGray,rounded corners=2pt,
    dashed,line width=.82pt,inner xsep=1.6pt,inner ysep=1.6pt},
  fail/.style={
    draw=SLWarningRed,fill=white,rounded corners=2pt,
    dash dot,line width=.82pt,inner xsep=1.6pt,inner ysep=1.6pt},
  statusband/.style={
    draw=SLTextGray,fill=white,rounded corners=2pt,
    line width=.80pt,inner xsep=1.6pt,inner ysep=1.6pt},
  ptitle/.style={
    font=\fontsize{10.2pt}{12.2pt}\selectfont\bfseries,
    inner sep=0pt},
  main/.style={->,draw=SLDeepBlue,line width=.90pt},
  success/.style={->,draw=SLSelfCheckTeal,line width=.90pt},
  budgetedge/.style={->,draw=SLAccentOrange,line width=.82pt,dashed},
  failure/.style={->,draw=SLWarningRed,line width=.82pt,dash dot},
  >={Stealth[length=2.1mm,width=1.5mm]},
  alt={对象—关系—结论：图由内容驱动的上下两面板组成。面板a把前置及空文档门、非空初始化与同步候选、可行门及原子提交和局部双证书排成一行，下一行分别承接输入失败、数值失败、局部预算停止和唯一LocalOK；非空初始化含均匀责任度、先验加平均词数及有限初始ELBO，只有completed或converged进入LocalOK。面板b先逐文档收齐全部LocalOK，再沿第二行从右向左完成全语料计数、严格正归一、受保护Newton、完整提交和外层双证书；下一行分开外层预算候选、收敛候选和F/X非成功状态，最终仅以feasible且状态为converged或budget_stop的记录参加跨启动选择，输出保留计数、诊断和原状态。}]

% ==========================================================================
% (a) Content-driven main row.
\matrix (amain) [
  column sep=3mm,
  inner sep=0pt,
  ampersand replacement=\&
] {
  \node[localgate,text width=40mm,minimum height=29mm] (a01)
    {\textbf{A0+A1 前置／空文档门}\\
     $\mathbf w_m;\ \alpha>0$\\
     $\varphi_{k\cdot}\in\operatorname{ri}(\Delta^{V-1})$\\
     $T_{\rm loc};\quad
     \varepsilon_{\rm loc}^{\mathcal L},
     \varepsilon_{\rm loc}^{\gamma}$\\
     词表、支持、预算、容差合法？\\
     $N_m=0$：$\mathtt{completed}\to L$\\
     $N_m>0\to$ A2；非法 $\to I$};
  \&
  \node[localbox,text width=50mm,minimum height=29mm] (a2)
    {\textbf{A2 非空初始化／同步候选（$R_A$ 入口）}\\
     $\eta_{mnk}^{(0)}=1/K;\quad
     \gamma_{mk}^{(0)}=\alpha_k+N_m/K$\\
     $\mathcal L_m^{(0)}\in\mathbb R$（有限）\\
     $a_{mnk}=\log\varphi_{k,w_{mn}}+\psi(\gamma_{mk})$\\
     $\hphantom{a_{mnk}=}-\psi(\sum_\ell\gamma_{m\ell})$\\
     $\eta^+_{mnk}=e^{a_{mnk}}\big/\sum_\ell e^{a_{mn\ell}}$\\
     （LSE；同步全部 $n$）\\
     $\gamma^+_{mk}=\alpha_k+\sum_n\eta^+_{mnk}$};
  \&
  \node[localcert,text width=52mm,minimum height=29mm] (a35)
    {\textbf{A3--A5 可行门／提交／局部双证书}\\
     $F_{\rm loc}$：$\eta^+\in\operatorname{ri}(\Delta),\quad
     \gamma^+>0$\\
     $\mathcal L_m^+\in\mathbb R,\quad
     \mathcal L_m^+\ge\mathcal L_m$\\
     通过才提交；$s_{\rm done}\mathrel{+}=1,\quad
     p_{\rm done}\mathrel{+}=N_m$\\
     $\delta_{\mathcal L,m}=
     |\mathcal L_m^+-\mathcal L_m|/(1+|\mathcal L_m|)$\\
     $\delta_{\gamma,m}=
     \max_k|\gamma^+_{mk}-\gamma_{mk}|/(1+|\gamma_{mk}|)$\\
     $C_{\rm loc}$：$\delta_{\mathcal L,m}\le
     \varepsilon_{\rm loc}^{\mathcal L}$ 且\\
     $\delta_{\gamma,m}\le\varepsilon_{\rm loc}^{\gamma}$？
     否且 $s_{\rm done}<T_{\rm loc}\to R_A$};
  \\
};

\node[ptitle,text=SLMidBlue,above=1mm of amain] (atitle)
  {(a) 单文档局部 VI（$\mathrm{LL}=\mathtt{last\ legal}$）};

% (a) Four compact direct status cells; true boundary gap is 3 mm.
\matrix (astatus) [
  below=3mm of amain,
  column sep=3mm,
  inner sep=0pt,
  ampersand replacement=\&
] {
  \node[fail,text width=31mm,minimum height=12mm] (ainvalid)
    {\textbf{$I:\ \mathtt{invalid\_input}$}\\
     $\mathtt{support\_failure}$\\
     存 $\mathrm{LL}/\mathtt{diag}$；\textbf{停／禁 C/M}};
  \&
  \node[fail,text width=30mm,minimum height=12mm] (anum)
    {\textbf{$\mathtt{numerical\_failure}$}\\
     回滚 $\mathrm{LL}$；gate／位置\\
     存 diag；\textbf{停止；禁 C/M}};
  \&
  \node[budget,text width=35mm,minimum height=12mm] (alocbudget)
    {\textbf{STOP：$\mathtt{budget\_stop}$}\\
     $\neg C_{\rm loc}$；$s_{\rm done}=T_{\rm loc}$\\
     存 $\mathrm{LL}$/计数/diag；$\nexists L$；\textbf{禁 C/M}};
  \&
  \node[ok,text width=42mm,minimum height=12mm] (aok)
    {\textbf{$L:\ \mathrm{LocalOK}_m$}\\
     $(\gamma_m,\eta_m,\mathcal L_m;\ \mathtt{counts},\mathtt{diag})$\\
     $\mathtt{completed}/\mathtt{converged}$};
  \\
};

% Only adjacent main arrows; the fan-out occupies the empty 3 mm gutter.
\draw[main] (a01) -- (a2);
\draw[main] (a2) -- (a35);
\draw[failure] (a01.south west) -- (ainvalid.north);
\draw[failure] (a35.south west) -- (anum.north);
\draw[budgetedge] (a35.south) -- (alocbudget.north);
\draw[success] (a35.south east) -- (aok.north);

% ==========================================================================
% (b) Title is content-positioned below panel (a); panel backgrounds are fit
%     only after every node has its final natural bounding box.
\node[ptitle,text=SLDeepBlue,below=3mm of astatus] (btitle)
  {(b) 语料级外层 VEM};

\matrix (btop) [
  below=1mm of btitle,
  column sep=3mm,
  inner sep=0pt,
  ampersand replacement=\&
] {
  \node[box,text width=52mm,minimum height=19mm] (b0)
    {\textbf{B0 全局门／多启动}\\
     $\mathcal D,K;\ \alpha,\varphi>0;\quad T_{\rm loc},T_{\rm out}$\\
     $\varepsilon_{\rm out}^{\mathcal L,\Theta};\quad
     s=1{:}R_{\rm start};\ J_{\rm ls}$\\
     随机源失败：$\mathtt{random\_source\_failure}$\\
     其余前置门否：$\mathtt{invalid\_input}$\\
     通过 $\to r_{\rm done}=0$};
  \&
  \node[box,text width=54mm,minimum height=19mm] (b1)
    {\textbf{B1 全部文档（$R$ 入口）}\\
     对 $m=1,\ldots,M$ 调用面板 (a)\\
     收集 $(\mathrm{LocalOK}_m,\mathtt{status}_m,
     \mathcal L_m,\mathtt{diag}_m)$\\
     局部 $\mathtt{budget\_stop}\to X$\\
     局部 failure $\to F$};
  \&
  \node[gate,text width=28mm,minimum height=19mm] (b2)
    {\textbf{B2 全收齐门}\\
     $\forall m:\ \mathrm{LocalOK}_m$ 到齐？\\
     是 $\to$ B3\\
     否 $\to F/X$（按原 status）};
  \\
};

% Right-to-left second row: B3 -> B4 -> B5/B6.
\matrix (bmain) [
  below=3mm of btop,
  column sep=3mm,
  inner sep=0pt,
  ampersand replacement=\&
] {
  \node[cert,text width=56mm,minimum height=29mm] (b56)
     {\textbf{B5--B6 提交／证书（$R:\ r<T_{\rm out}$）}\\
     $\Theta^+=(\alpha^+,\varphi^+,
     \{\gamma_m^+,\eta_m^+\}_{m=1}^{M},\mathcal L^+)$\\
     可行、有限且不下降才提交；
     $r_{\rm done}\mathrel{+}=1$\\
     $\delta_{\mathcal L}=
     |\mathcal L^+-\mathcal L|/(1+|\mathcal L|)$\\
     $\delta_\alpha=
     \max_k|\alpha_k^+-\alpha_k|/(1+|\alpha_k|)$\\
     $\delta_\varphi=\max_{k,v}
     |\varphi^+_{kv}-\varphi_{kv}|$\\
     $\delta_\Theta=\max\{\delta_\alpha,\delta_\varphi\}$\\
     $C_{\rm out}$：
     $\delta_{\mathcal L}\le\varepsilon_{\rm out}^{\mathcal L}
     \land\delta_\Theta\le\varepsilon_{\rm out}^{\Theta}$？};
  \&
  \node[cert,text width=40mm,minimum height=29mm,inner ysep=2.4pt] (b4)
    {\textbf{B4 受保护 Newton}\\
     \textbf{（$N$ 入口）}\\
     $H(\alpha)\Delta=g(\alpha)$\\
     $\alpha^+=\alpha-\rho\Delta$\\
     $\rho\in\{1,\tfrac12,\tfrac14,\ldots\}$；
     $j\le J_{\rm ls}$\\
     $G_N$：$\min_k\alpha_k^+>0$\\
     且新 ELBO 有限、不下降\\
     拒绝且 $j<J_{\rm ls}$：
     $\rho\leftarrow\rho/2\to N$\\
     $j=J_{\rm ls}$：
     $\mathtt{line\_search\_failed}\to F$};
  \&
  \node[box,text width=42mm,minimum height=29mm] (b3)
    {\textbf{B3 全语料计数 $\to\varphi^+$}\\
     $C^+_{kv}=\sum_{m=1}^{M}\sum_{n=1}^{N_m}$\\
     $\eta_{mnk}\mathbf1\{w_{mn}=v\}$\\
     $G_M$：计数有限；$C^+_{k\cdot}>0$\\
     且每个 $C^+_{kv}>0$\\
     $\varphi^+_{kv}=C^+_{kv}/C^+_{k\cdot}>0$\\
     $\sum_v\varphi^+_{kv}=1$\\
     失败 $\to F:\mathtt{numerical\_failure}$};
  \\
};

% Compact two-row terminal table.  Its top row is addressable by all exits;
% the full-width selection row is positioned from the completed top bbox.
\matrix (bterminal) [
  below=3mm of bmain,
  column sep=3mm,
  inner sep=0pt,
  ampersand replacement=\&
] {
  \node[budget,text width=48mm,minimum height=12mm] (bbudget)
    {\textbf{STOP：$\mathtt{budget\_stop}$（未收敛候选）}\\
     $\neg C_{\rm out},\ r=T_{\rm out}$（本启动）\\
     可行；存 $\mathrm{LL}$/计数/diag（原 status）};
  \&
  \node[ok,text width=36mm,minimum height=12mm] (bconv)
    {$\checkmark$ \textbf{$\mathtt{converged}$ 候选}\\
     $C_{\rm out}$；$\mathtt{feasible}=\mathtt{true}$\\
     存 $\mathtt{record}[s]$（原 status）};
  \&
  \node[statusband,text width=58mm,minimum height=12mm] (bfail)
    {\textbf{$F/X$：非成功；禁 C/M；$\notin\mathcal S_{\rm acc}$}\\
     $F$：保留入口的原失败 status\\
     $X$：局部 $\mathtt{budget\_stop}$（原 status）\\
     $\mathtt{feasible}=\mathtt{false}$；仅存 status／$\mathrm{LL}$/diag};
  \\
};

\node[box,text width=148mm,minimum height=15mm,fill=SLSoftTeal,
  font=\fontsize{9.6pt}{9.9pt}\selectfont,
  below=2mm of bterminal] (bselect)
  {\textbf{跨启动选择／输出}\quad $q_s:=\mathtt{record}[s]$\\
   $\mathcal S_{\rm acc}=\{s:q_s.\mathtt{feasible}=\mathtt{true},\quad
   q_s.\mathtt{status}\in\{\mathtt{converged},\mathtt{budget\_stop}\}\}$\\
   空集：仅返回全部 records；非空：按最终 ELBO／冻结并列规则返回
   $(\widehat\alpha,\widehat\varphi,
   \{\widehat\gamma_m,\widehat\eta_m\},\mathcal L)$\\
   $s_{\rm done},r_{\rm done},m_{\rm done},k_{\rm done},j_{\rm done}$；
   diagnostics；原 status：
   \textcolor{SLSelfCheckTeal}{$\checkmark\mathtt{converged}$}／
   \textcolor{SLAccentOrange}{STOP $\mathtt{budget\_stop}$（未收敛）}};

% Main edges are only between adjacent nodes. Cross-row edges stay in the
% 3 mm boundary gutters and preserve left-to-right endpoint ordering.
\draw[main] (b0) -- (b1);
\draw[main] (b1) -- (b2);
\draw[main] (b2.south) -- (b3.north);
\draw[main] (b3) -- (b4);
\draw[main] (b4) -- (b56);

% B3 and B4 each have a real, short failure edge to the right-hand F/X cell.
\draw[failure] (b3.south) -- (bfail.north east);
\draw[failure] (b4.south east) -- (bfail.north west);
\draw[success] (b56.south east) -- (bconv.north);
\draw[budgetedge] (b56.south west) -- (bbudget.north);

\draw[budgetedge] (bbudget.south) -- (bselect.north west);
\draw[success] (bconv.south) -- (bselect.north);

% Content-derived backgrounds are created last, after every real node bbox.
\begin{scope}[on background layer]
  \node[
    draw=SLMidBlue,fill=SLSoftTeal,fill opacity=.24,
    line width=.80pt,rounded corners=4pt,
    fit=(atitle)(amain)(astatus),
    inner xsep=2mm,inner ysep=1mm] (apanel) {};
  \node[
    draw=SLDeepBlue,fill=SLSoftBlue,fill opacity=.22,
    line width=.80pt,rounded corners=4pt,
    fit=(btitle)(btop)(bmain)(bterminal)(bselect),
    inner xsep=2mm,inner ysep=1mm] (bpanel) {};
\end{scope}

\end{tikzpicture}%
}
\captionsetup{width=.97\linewidth}
\caption{点参数 LDA 两层变分 EM：仅 $\mathtt{completed}/\mathtt{converged}$
形成 $\mathrm{LocalOK}_m$，全部 $m$ 到齐后才执行
$C\to\varphi\to$ 受保护 Newton；可行的 $\mathtt{budget\_stop}$
保留原 status，作为未收敛候选进入 $\mathcal S_{\rm acc}$，
其余非成功状态止于记录。}
\label{fig:V5-C06-variational-updates}
\end{figure}

\clearpage
\typeout{SLFIG-HARNESS-END FIG-P694-01}

\typeout{SLFIG-HARNESS-BEGIN FIG-P722-01}
% v2.3.1 figure UID: FIG-P722-01
% v2.6.0 reverse redraw for FIG-P722-01: preserve the exact ten-node/ten-edge sparse-power topology while routing feedback through real exterior whitespace.
\tikzset{slfig-FIG-P722-01/.style={font=\fontsize{9.2pt}{11.0pt}\selectfont,every path/.append style={line cap=round,line join=round}}}
\pgfkeysifdefined{/tikz/alt}{}{\tikzset{alt/.code={}}}
\begin{figure}[H]
\centering
\begin{tikzpicture}[slfig-FIG-P722-01,
  every node/.style={font=\fontsize{9.2pt}{11.0pt}\selectfont,align=center,outer sep=0pt},
  stepbox/.style={draw=SLMidBlue,fill=SLSoftBlue,rounded corners=2pt,minimum height=12.5mm,text width=29mm,inner xsep=2.0mm,inner ysep=1.4mm},
  inputbox/.style={stepbox,draw=SLRuleGray,fill=SLSoftGray},
  check/.style={diamond,draw=SLDeepBlue,fill=white,aspect=2.20,inner sep=1.4pt,text width=20mm},
  successbox/.style={draw=SLDeepBlue,fill=SLDeepBlue,text=white,rounded corners=2pt,text width=35mm,minimum height=12mm,inner xsep=2.0mm,inner ysep=1.4mm},
  budgetbox/.style={draw=SLRuleGray,fill=SLSoftGray,rounded corners=2pt,text width=35mm,minimum height=12mm,inner xsep=2.0mm,inner ysep=1.4mm},
  certbox/.style={draw=SLMidBlue,fill=white,rounded corners=2pt,text width=37mm,minimum height=11.5mm,inner xsep=2.0mm,inner ysep=1.3mm},
  main/.style={->,draw=SLDeepBlue,line width=.95pt},
  loop/.style={->,draw=SLRuleGray,line width=.72pt,densely dashed},
  edgelabel/.style={font=\fontsize{9.2pt}{11.0pt}\selectfont,inner sep=0pt},
  >={Stealth[length=2.2mm,width=1.55mm]},
  alt={对象—关系—结论：七步稀疏幂法依次检查输入、计算悬挂质量、稀疏乘法、加入阻尼跳转、归一化与非负核对、计算相邻迭代差，再由判定节点分成收敛、继续或预算停止三路；继续回路只返回悬挂质量步骤。}]
\node[inputbox] (s1) at (-6.00,2.55) {1 输入检查\\维数、$d,v,r^{(0)}$};
\node[stepbox] (s2) at (-2.00,2.55) {2 悬挂质量\\$m_t=a^{\mathsf T}r^{(t)}$};
\node[stepbox] (s3) at (2.00,2.55) {3 稀疏乘法\\$q_t=Mr^{(t)}$};
\node[stepbox] (s4) at (6.00,2.55) {4 阻尼／跳转\\$r^+=d(q_t+m_tv)+(1-d)v$};
\node[stepbox] (s5) at (6.00,0.15) {5 归一化／非负核对\\$r^+\in\Delta^{n-1}$};
\node[stepbox] (s6) at (2.00,0.15) {6 相邻差\\$\delta_t=\lVert r^+-r^{(t)}\rVert_1$};
\node[check] (s7) at (-2.20,0.15) {7 停止条件\\或预算？};
\node[successbox] (conv) at (-5.00,-2.60) {converged\\输出$r^+$};
\node[budgetbox] (budget) at (1.20,-2.60) {budget\_stop\\报告$r^+,\delta_t$};
\node[certbox] (cert) at (-5.00,-4.25) {$\delta_t/(1-d)\le\varepsilon$\\固定点误差证书};
\draw[main] (s1)--(s2); \draw[main] (s2)--(s3); \draw[main] (s3)--(s4); \draw[main] (s4)--(s5); \draw[main] (s5)--(s6); \draw[main] (s6)--(s7);
\draw[main] (s7.south west)--(conv.north east) node[edgelabel,pos=.52,left=3.0mm]{收敛};
\draw[main] (s7.south east)--(budget.north west) node[edgelabel,pos=.52,right=3.0mm]{预算到};
\draw[main] (conv.south)--(cert.north);
\draw[loop] (s7.west)--(-8.25,.15)--(-8.25,4.30)--(-2.00,4.30)--(s2.north)
  node[edgelabel,pos=.40,left=2.0mm]{continue};
\end{tikzpicture}
\captionsetup{width=.94\linewidth}
\caption{PageRank稀疏幂法依次执行输入检查、悬挂质量回填、稀疏乘法与阻尼更新，并以$\delta_t=\lVert r^{(t+1)}-r^{(t)}\rVert_1\le(1-d)\varepsilon$作为$\ell_1$误差证书；不满足时继续迭代或按预算停止。}
\label{fig:V5-C07-power-flow}
\end{figure}

\clearpage
\typeout{SLFIG-HARNESS-END FIG-P722-01}

\typeout{SLFIG-HARNESS-BEGIN FIG-P756-01}
\input{../../绘图源码/第05册_采样方法主题模型与图排序/V5-C08/full_course_synthesis_map.tex}
\clearpage
\typeout{SLFIG-HARNESS-END FIG-P756-01}

\end{document}
